% 四元素底集 {1,2,3,4} 的子集格；框内上行为子集 S，下行为自身贡献 h(S)。
% 教学设定：单项各2，{1,3} 额外3，{2,4} 额外-1，其余0。蓝色为 {1,3} 及其全部子集，累计 g({1,3})=0+2+2+3=7。
\begin{tikzpicture}[x=1mm,y=1mm,line width=.6pt,font=\figurefont\fontsize{8}{10}\selectfont,
  s/.style={draw=FigureStroke,fill=white,rounded corners=1.8pt,align=center,inner sep=1.2pt,minimum width=12.5mm,minimum height=8mm,line width=.8pt}]
  \coordinate (cne) at (0.0,0.0);
  \coordinate (cn1) at (-25.5,15.0);
  \coordinate (cn2) at (-8.5,15.0);
  \coordinate (cn3) at (8.5,15.0);
  \coordinate (cn4) at (25.5,15.0);
  \coordinate (cn12) at (-42.5,30.0);
  \coordinate (cn13) at (-25.5,30.0);
  \coordinate (cn14) at (-8.5,30.0);
  \coordinate (cn23) at (8.5,30.0);
  \coordinate (cn24) at (25.5,30.0);
  \coordinate (cn34) at (42.5,30.0);
  \coordinate (cn123) at (-25.5,45.0);
  \coordinate (cn124) at (-8.5,45.0);
  \coordinate (cn134) at (8.5,45.0);
  \coordinate (cn234) at (25.5,45.0);
  \coordinate (cn1234) at (0.0,60.0);
  \draw[FigureGrid,line width=.6pt] (cne)--(cn2);
  \draw[FigureGrid,line width=.6pt] (cne)--(cn4);
  \draw[FigureGrid,line width=.6pt] (cn1)--(cn12);
  \draw[FigureGrid,line width=.6pt] (cn1)--(cn14);
  \draw[FigureGrid,line width=.6pt] (cn2)--(cn12);
  \draw[FigureGrid,line width=.6pt] (cn2)--(cn23);
  \draw[FigureGrid,line width=.6pt] (cn2)--(cn24);
  \draw[FigureGrid,line width=.6pt] (cn3)--(cn23);
  \draw[FigureGrid,line width=.6pt] (cn3)--(cn34);
  \draw[FigureGrid,line width=.6pt] (cn4)--(cn14);
  \draw[FigureGrid,line width=.6pt] (cn4)--(cn24);
  \draw[FigureGrid,line width=.6pt] (cn4)--(cn34);
  \draw[FigureGrid,line width=.6pt] (cn12)--(cn123);
  \draw[FigureGrid,line width=.6pt] (cn12)--(cn124);
  \draw[FigureGrid,line width=.6pt] (cn13)--(cn123);
  \draw[FigureGrid,line width=.6pt] (cn13)--(cn134);
  \draw[FigureGrid,line width=.6pt] (cn14)--(cn124);
  \draw[FigureGrid,line width=.6pt] (cn14)--(cn134);
  \draw[FigureGrid,line width=.6pt] (cn23)--(cn123);
  \draw[FigureGrid,line width=.6pt] (cn23)--(cn234);
  \draw[FigureGrid,line width=.6pt] (cn24)--(cn124);
  \draw[FigureGrid,line width=.6pt] (cn24)--(cn234);
  \draw[FigureGrid,line width=.6pt] (cn34)--(cn134);
  \draw[FigureGrid,line width=.6pt] (cn34)--(cn234);
  \draw[FigureGrid,line width=.6pt] (cn123)--(cn1234);
  \draw[FigureGrid,line width=.6pt] (cn124)--(cn1234);
  \draw[FigureGrid,line width=.6pt] (cn134)--(cn1234);
  \draw[FigureGrid,line width=.6pt] (cn234)--(cn1234);
  \draw[FigureBlue,line width=1.4pt] (cne)--(cn1);
  \draw[FigureBlue,line width=1.4pt] (cne)--(cn3);
  \draw[FigureBlue,line width=1.4pt] (cn1)--(cn13);
  \draw[FigureBlue,line width=1.4pt] (cn3)--(cn13);
  \node[s,fill=FigureBlueFill] (ne) at (0.0,0.0) {$\varnothing$\\$h=0$};
  \node[s,fill=FigureBlueFill] (n1) at (-25.5,15.0) {$\{1\}$\\$h=2$};
  \node[s,fill=white] (n2) at (-8.5,15.0) {$\{2\}$\\$h=2$};
  \node[s,fill=FigureBlueFill] (n3) at (8.5,15.0) {$\{3\}$\\$h=2$};
  \node[s,fill=white] (n4) at (25.5,15.0) {$\{4\}$\\$h=2$};
  \node[s,fill=white] (n12) at (-42.5,30.0) {$\{1,2\}$\\$h=0$};
  \node[s,fill=FigureBlueFill] (n13) at (-25.5,30.0) {$\{1,3\}$\\$h=3$};
  \node[s,fill=white] (n14) at (-8.5,30.0) {$\{1,4\}$\\$h=0$};
  \node[s,fill=white] (n23) at (8.5,30.0) {$\{2,3\}$\\$h=0$};
  \node[s,fill=FigureYellowFill] (n24) at (25.5,30.0) {$\{2,4\}$\\$h=-1$};
  \node[s,fill=white] (n34) at (42.5,30.0) {$\{3,4\}$\\$h=0$};
  \node[s,fill=white] (n123) at (-25.5,45.0) {$\{1,2,3\}$\\$h=0$};
  \node[s,fill=white] (n124) at (-8.5,45.0) {$\{1,2,4\}$\\$h=0$};
  \node[s,fill=white] (n134) at (8.5,45.0) {$\{1,3,4\}$\\$h=0$};
  \node[s,fill=white] (n234) at (25.5,45.0) {$\{2,3,4\}$\\$h=0$};
  \node[s,fill=white] (n1234) at (0.0,60.0) {$\{1,2,3,4\}$\\$h=0$};
\end{tikzpicture}
